\section{The score-normalisation artefact}
\label{sec:artefact}

The original's headline empirical conclusion is that ESG-washing is largely
absent from the reports it studies: \emph{``The ESGSI predominantly shows
negative values, suggesting a general absence of greenwashing practices among
the entities studied''} \citep[p.~6]{lagasio2024}. This section argues that the
conclusion is an artefact of two methodological choices, neither of which the
paper states in a form a reader could act on --- the \emph{score normalisation}
applied cross-sectionally to the two components at index assembly, which the paper
specifies twice and inconsistently, and the specification of the substance
component (\S\ref{sec:index}), whose aggregation rule it never gives at all
--- rather than a property of the disclosure.

The argument has three parts, and they are deliberately kept separate because
they do not carry equal weight. The first establishes, from the original's own
published tables, \emph{which} score normalisation produced its numbers; this
rests on arithmetic that requires nothing of our pipeline. The second shows
analytically why that choice determines the sign of the reported mean; it too
requires nothing of our pipeline, and nothing of the original beyond the two
maps themselves. The third demonstrates the consequence on our corpus. The
first two are jointly sufficient for the section's claim and neither uses our
data; the third is corroboration, and we mark it as such.

\subsection{Which score normalisation produced the published numbers}
\label{sec:artefact-which}

Section 3.3 of the original states its score normalisation in prose:

\begin{quote}
``To ensure comparability, both sentiment and sustainability scores are
normalized to a uniform scale using min-max normalization. This normalization
technique is widely used in text analysis and machine learning to address
potential biases due to variations in document length and linguistic style
(Sebastiani, 2002).'' \citep[p.~4]{lagasio2024}
\end{quote}

\noindent The formula printed a few lines below, in the same subsection, is a
difference of $z$-scores:
\begin{equation}
\label{eq:artefact-printed}
\ESGSI_i \;=\; \frac{\SEN_i - \overline{\SEN}}{\sigma_{\SEN}}
\;-\; \frac{\SUS_i - \overline{\SUS}}{\sigma_{\SUS}}
\end{equation}
with the six symbols defined immediately beneath it as the component scores,
their sample means and their sample standard deviations
\citep[p.~4]{lagasio2024}. The paper never reconciles the two; we flagged the
contradiction at the point of definition (\S\ref{sec:index}) and resolve it here.

Four lines of evidence, all drawn from the original's own Table~2 ($n = 749$),
settle it against the printed formula.

\emph{Support.} The table notes state verbatim that the \emph{``Normalized
Sustainability Score: Scores range from 0 (low sustainability) to 1 (high
sustainability)''} and that the \emph{``Normalized Sentiment Score \dots
[is] normalized between 0 and 1''} \citep[p.~6]{lagasio2024}; the note to
Figure~1 repeats that the plotted component scores run \emph{``over a
continuous range from 0 to 1''} \citep[p.~5]{lagasio2024}. The reported
minima, $0.0546$ for the sustainability component and $0.0000$ for the
sentiment component, are both non-negative. A $z$-scored variable has mean
zero and positive dispersion, so its minimum is necessarily negative. Neither
component can have been $z$-scored.

\emph{Dispersion.} Table~2 gives standard deviations of $0.1321$ for the
sustainability component and $0.0853$ for the sentiment component. A $z$-score
has unit variance by construction. Two variables with standard deviations of
$0.1321$ and $0.0853$ have not been $z$-scored.

\emph{Location.} Table~2 reports mean $\ESGSI = -0.2179$ with standard
deviation $0.1614$. The sentiment component has mean $0.3727$ and the
sustainability component mean $0.5906$, and
\begin{equation}
\label{eq:artefact-identity}
0.3727 - 0.5906 \;=\; -0.2179 ,
\end{equation}
which is the reported mean exactly. Under Equation~\eqref{eq:artefact-printed}
the sample mean of $\ESGSI$ is $0$ by construction, since it is a difference of
two mean-zero variables. It is not $0$.

\emph{Quantiles.} Table~2 also reports quartiles. For the index they are
$-0.3245$, $-0.2177$ and $-0.1054$ \citep[p.~6]{lagasio2024}: the third
quartile lies below the labelling threshold, so at least three quarters of the
$749$ reports fall at or below $-0.1054$, and the share the original's own
instrument would have labelled \emph{``Potential ESGwashing''} is bounded above
at roughly one report in four --- at most $188$ of $749$ on the most
conservative of the standard quantile conventions. The original never reports
that share (\S\ref{sec:original}); its own table bounds it. The bound is hard
to reconcile with Equation~\eqref{eq:artefact-printed}. Under the printed
formula the index has mean zero by construction, so for only a quarter of the
sample to exceed zero the index would have to be pronouncedly right-skewed,
with its mean at or above its own third quartile. Nothing in Table~2 supplies
such a skew: the reported quartiles are close to evenly spaced about the median
for the sustainability component ($0.5020 / 0.5876 / 0.6732$, gaps of $0.0856$
and $0.0856$), for the sentiment component ($0.3188 / 0.3668 / 0.4169$, gaps of
$0.0480$ and $0.0501$), and for the reported index itself, whose mean
($-0.2179$) and median ($-0.2177$) differ by $0.0002$.

The identity in Equation~\eqref{eq:artefact-identity} is not confined to the
overall row. Subtracting the reported ``Normalized Sustainability Score'' from
the reported ``Normalized Sentiment Score'' reproduces the reported $\ESGSI$ in
every one of the twenty-eight rows the original publishes across Tables~2--6:
exactly in twenty-four of them, and to within one unit in the fourth decimal
place --- the resolution at which the components themselves are printed --- in
the remaining four (Table~\ref{tab:artefact-rows}). No row was selected or
omitted; the table below is the complete published set.

\begin{table}[htbp]
\centering
\caption{Reconstruction of the original's reported $\ESGSI$ as the difference
of its two reported components. Every row published in \citet{lagasio2024},
Tables~2--6, is included. Four rows --- Consumer Staples, the Table~4 total,
and the second and fourth ESG quartiles --- differ by one unit in the fourth
decimal place, which is the resolution at which the two components are printed.
All figures are read from the original.}
\label{tab:artefact-rows}
\begin{tabular}{lrrrr}
\toprule
Row & Norm.\ $\SEN$ & Norm.\ $\SUS$ & Difference & Reported $\ESGSI$ \\
\midrule
\multicolumn{5}{l}{\emph{Table 2, p.~6 --- full sample}} \\
Overall                & 0.3727 & 0.5906 & $-0.2179$ & $-0.2179$ \\
\midrule
\multicolumn{5}{l}{\emph{Table 3, p.~7 --- by region}} \\
Europe                 & 0.3477 & 0.6144 & $-0.2667$ & $-0.2667$ \\
North America          & 0.3917 & 0.5731 & $-0.1814$ & $-0.1814$ \\
South America          & 0.3229 & 0.6622 & $-0.3393$ & $-0.3393$ \\
Other                  & 0.3665 & 0.5760 & $-0.2095$ & $-0.2095$ \\
Total                  & 0.3572 & 0.6064 & $-0.2492$ & $-0.2492$ \\
\midrule
\multicolumn{5}{l}{\emph{Table 4, p.~7 --- by GICS sector}} \\
Communication Services & 0.3637 & 0.5909 & $-0.2272$ & $-0.2272$ \\
Consumer Discretionary & 0.4087 & 0.5889 & $-0.1802$ & $-0.1802$ \\
Consumer Staples       & 0.4048 & 0.6138 & $-0.2090$ & $-0.2089$ \\
Energy                 & 0.3478 & 0.6595 & $-0.3117$ & $-0.3117$ \\
Financials             & 0.3460 & 0.5944 & $-0.2484$ & $-0.2484$ \\
Health Care            & 0.3832 & 0.5775 & $-0.1943$ & $-0.1943$ \\
Industrials            & 0.3850 & 0.5743 & $-0.1893$ & $-0.1893$ \\
Information Technology & 0.3788 & 0.5547 & $-0.1759$ & $-0.1759$ \\
Materials              & 0.3584 & 0.6187 & $-0.2603$ & $-0.2603$ \\
Real Estate            & 0.3634 & 0.5672 & $-0.2038$ & $-0.2038$ \\
Utilities              & 0.3483 & 0.5873 & $-0.2390$ & $-0.2390$ \\
Total                  & 0.3716 & 0.5934 & $-0.2218$ & $-0.2217$ \\
\midrule
\multicolumn{5}{l}{\emph{Table 5, p.~8 --- by total-assets quartile}} \\
1st quartile           & 0.3853 & 0.5688 & $-0.1835$ & $-0.1835$ \\
2nd quartile           & 0.3829 & 0.5878 & $-0.2049$ & $-0.2049$ \\
3rd quartile           & 0.3620 & 0.5949 & $-0.2329$ & $-0.2329$ \\
4th quartile           & 0.3595 & 0.6093 & $-0.2498$ & $-0.2498$ \\
Total                  & 0.3724 & 0.5902 & $-0.2178$ & $-0.2178$ \\
\midrule
\multicolumn{5}{l}{\emph{Table 6, p.~9 --- by ESG-score quartile}} \\
1st quartile           & 0.3857 & 0.5650 & $-0.1793$ & $-0.1793$ \\
2nd quartile           & 0.3772 & 0.5898 & $-0.2126$ & $-0.2127$ \\
3rd quartile           & 0.3751 & 0.5855 & $-0.2104$ & $-0.2104$ \\
4th quartile           & 0.3529 & 0.6223 & $-0.2694$ & $-0.2695$ \\
Total                  & 0.3727 & 0.5907 & $-0.2180$ & $-0.2180$ \\
\bottomrule
\end{tabular}
\end{table}

One passage in the original describes its index in terms that fit the printed
formula rather than the tables. Introducing Figure~1, the paper states that the
distribution \emph{``is bell-shaped and centered around an ESGSI of zero''} and
that \emph{``a symmetric dispersion pattern around the central value indicates
an equivalent propensity for documents to exhibit both potential greenwashing
behaviors and genuine sustainability representation''} \citep[p.~6]{lagasio2024}.
That description is not consistent with Table~2 on the same page --- mean
$-0.2179$, median $-0.2177$, third quartile $-0.1054$, all below zero --- nor
with the paper's own summary of that table, two paragraphs later, that the index
\emph{``predominantly shows negative values''}. We do not attempt to explain the
discrepancy, and we do not treat prose as evidence about which map was applied:
the determination above rests on the numeric tables, which are also what every
downstream analysis in the original --- Tables~3--7 and the ANOVA --- is
computed from.

The estimated index is therefore the difference of the two reported components
on their $[0,1]$ scale, the min--max rescaling of the original's own prose, and
not the $z$-score difference the paper prints. What the four arguments establish
jointly is the negative, and it is the negative that the section needs: the
published numbers were not produced by Equation~\eqref{eq:artefact-printed}.
Nothing in that conclusion requires the rescaling to have been min--max in every
operational detail. It requires only that the components be bounded and that the
index be their plain difference, and the original's tables show both directly.

This argument uses only numbers the original published. It would stand unchanged
if our corpus, our vocabulary and our pipeline did not exist, and it is not
revisited, reinforced or qualified anywhere later in this paper. The choice is
established here, and only here.

\subsection{Why the choice determines the sign of the result}
\label{sec:artefact-why}

Write $\Z(\cdot)$ for the $z$-score map and $M(\cdot)$ for the min--max map
applied cross-sectionally to a finished component score. Both are affine and
increasing, so within a component they are rank-equivalent; they are not
equivalent once two components are differenced, because they equalise the two
components' dispersions in different ways.

Under $\Z$, the sample mean of $\Z(\SEN) - \Z(\SUS)$ is $0$ identically,
whatever the two components contain: the estimator centres the index on its own
labelling threshold before any document is read. What remains free is only the
shape of the distribution about that centre, so whether a majority falls below
zero becomes a question about the skew of a mean-zero variable rather than about
disclosure practice --- and the original's own reported quartiles
(\S\ref{sec:artefact-which}) describe components that supply no such skew.
Under the printed formula, ``predominantly negative'' is neither available as a
result nor interpretable as one. Under $M$, the sample mean is
\begin{equation}
\label{eq:artefact-mm-mean}
\overline{M(\SEN)} - \overline{M(\SUS)},
\end{equation}
which the construction pins to no particular value; it is fixed by where each
component's mass sits inside its own observed range.
Equation~\eqref{eq:artefact-mm-mean} is a statement about the position of two
marginal distributions within their ranges, not about firm behaviour. The
original reports the two quantities directly --- substance $0.5906$, above the
midpoint of its own reported range, against sentiment $0.3727$, well below the
midpoint of its --- and their difference is the reported index mean, by
Equation~\eqref{eq:artefact-identity}. The original describes the same
asymmetry in words, calling its sustainability scores \emph{``skewed towards
higher values''} and its sentiment scores approximately bell-shaped
\citep[p.~6]{lagasio2024}. The sign of the mean is therefore a fact about where
two marginal distributions sit within their own ranges. It becomes a claim
about firm behaviour only through the zero threshold, which is definitional and
uncalibrated (\S\ref{sec:index}) and which inherits the asymmetry directly.

The paper's own justification for the choice does not survive inspection. The
sole stated defence is that min--max score normalisation addresses
\emph{``potential biases due to variations in document length and linguistic
style''} \citep[p.~4]{lagasio2024}, a claim repeated at the end of its
Section~3.2.3. Min--max score normalisation is a single monotone affine
rescaling applied across the cross-section, $(x - \min x)/(\max x - \min x)$.
It is rank-preserving within each variable, so whatever dependence on document
length is present in $\SUS_i$ before the rescaling is present, in the same rank
order, after it. A rank-preserving cross-sectional map cannot remove a
dependence that survives in the ranks. The operation that could --- vector
normalisation, applied to the document's term vector at the TF-IDF stage, whose
length-neutrality we demonstrate in \S\ref{sec:measurement} --- is credited with
nothing and is never named (\S\ref{sec:original}). The paper's only stated
protection against length bias is attributed to the operation that cannot
provide it, while the operation that can provide it goes unmentioned.

\subsection{The consequence, demonstrated on our corpus}
\label{sec:artefact-consequence}

We can now ask what the two choices do jointly. Because our corpus, ESG
vocabulary and sentiment dictionary all differ from the original's by design
(\S\ref{sec:scope}, \S\ref{sec:lexicon}), what follows is a demonstration of a
mechanism on an independent sample; it is not a replication of the original's
estimate and is not offered as further evidence for
\S\ref{sec:artefact-which}.

Table~\ref{tab:artefact-factorial} crosses the two score normalisations with
two substance specifications --- our density specification $\susdens$ and the
L2 vector-normalised TF-IDF specification $\suslag$ that follows the original's
weighting (\S\ref{sec:index}) --- over our 343 reports.

\begin{table}[htbp]
\centering
\caption{Mean index and number of reports flagged ($\ESGSI > 0$) under the
$2 \times 2$ factorial over score normalisation and substance specification,
$n = 343$. The two $z$-score means are zero by construction and are printed at
the same resolution as the others. Counts are quantities of our corpus and are
not estimates of any quantity of the original's sample.}
\label{tab:artefact-factorial}
\begin{tabular}{llrr}
\toprule
Score normalisation & Substance specification & Mean & Flagged \\
\midrule
$z$-score & $\susdens$ & $0.0000$ & $171/343$ \\
$z$-score & $\suslag$  & $0.0000$ & $159/343$ \\
min--max  & $\susdens$ & $-0.0004$ & $166/343$ \\
min--max  & $\suslag$  & $\mathbf{-0.2388}$ & $\mathbf{54/343}$ \\
\bottomrule
\end{tabular}
\end{table}

Only the joint configuration moves the mean appreciably away from zero; the
other three cells sit within $0.0004$ of it. Holding the substance
specification at $\suslag$ and switching the score normalisation to $z$-scores
returns the mean to zero exactly; holding min--max and switching the substance
specification to $\susdens$ returns it to within $0.0004$ of zero. The flagged
count falls from about half the corpus to $54/343$ only in the joint cell.

Three restrictions on how this may be read. First, our $-0.2388$ and the
original's $-0.2179$ differ by $0.0209$, and the correspondence between them is
consistency, not proof: a matching sign and order of magnitude obtained on a
different corpus with a different vocabulary corroborates
the mechanism and does not establish it, and coincidence is possible. The claim
we make is that these two choices jointly reproduce the statistic and are
jointly necessary for it \emph{on our corpus} --- not that we have shown what
the original ran, which \S\ref{sec:artefact-which} establishes on stronger and
independent grounds. Second, the flagged counts of $171/343$ and $54/343$
are quantities of our corpus alone and are not estimates of the original's
flagged share. The original never reports that share (\S\ref{sec:original}) and
nothing in this subsection recovers it; the only statement available about it is
the upper bound its own published quartiles impose, and that bound was obtained
in \S\ref{sec:artefact-which} without reference to anything here. Third, a
reader who rejects this subsection in its entirety is left with
\S\ref{sec:artefact-which} and \S\ref{sec:artefact-why} intact, and those are
sufficient for the section's claim.

\subsection{What follows}
\label{sec:artefact-scope}

The published conclusion does not survive the paper's own printed formula.
Under Equation~\eqref{eq:artefact-printed} the index is centred on its own
labelling threshold by construction, so ``predominantly negative'' is not a
substantive result the estimator can return; under the min--max rescaling
actually used, the sign of the mean is set by the position of two marginal
distributions within their own ranges and by the substance specification --- the
first specified twice and inconsistently, the second not specified at all. The
conclusion is thus a property of the index as assembled, and it is not stable
under alternatives that the paper's own text licenses --- which is the empirical
counterpart of the label instability we report in \S\ref{sec:robustness}.

This is a claim about the instrument, and its scope should not be overread. We
do not claim that ESG-washing is prevalent, nor that the firms in the
original's sample engage in it; establishing either would require an external
benchmark that neither the original nor this paper possesses
(\S\ref{sec:limitations}). What we claim is narrower and, we think, harder to
avoid: the reported absence of ESG-washing is a consequence of the score
normalisation applied and of how substance was specified, and neither choice
was disclosed in a form that would have let a reader see it.
